% !TEX root = ../../../main.tex

\subsection{المرحلة الثالثة: الترشيح الكامل (إضافة الرتب $n=7$ و $n=11$)}
وهي الخطوة المطلوبة في الوظيفة، حيث تم إدراج فرعين إضافيين على التوازي مع المرشح السابق:

\begin{enumerate}
    \item \textbf{مرشح التوافقية السابعة ($n=7$):} تم إدخال القيم المحسوبة نظرياً: $R=0.1815 \Omega$ و $L=0.00619 H$.
    
    \begin{figure}[H]
    \centering
    \includegraphics[width=0.75\linewidth]{08 - 2nd filter parameters.png} % Replace with actual file
    \caption{بارامترات مرشح التوافقية السابعة في ماتلاب}
    \label{fig:08}
    \vspace{-10pt}
\end{figure}
    
    \newpage
    \begin{figure}[H]
      \vspace{-30pt}
        \centering
        \includegraphics[angle=270,width=1\linewidth]{07 - Circuit with 2nd filter.png}
        \caption{شكل الدارة بعد إضافة المرشح الثاني ($n=7$)}
    \end{figure}
    
\newpage

عند إضافة المرشح الثاني ($n=7$)، نلاحظ من خلال أداة التحليل تلاشي المزيد من القمم، نتيجة ترشيح التوافقيات المقابلة لهذا التردد، مما أدى إلى تحسين إضافي في جودة القدرة:
\begin{figure}[H]
  \centering
  \includegraphics[width=0.9\textwidth]{09 - CurrentData with 2nd filter.png}
    \caption{\small نتائج تحليل \el{FFT} للتيار بعد اضافة المرشح الثاني - تظهر \el{THD} بنسبة $11.94\%$}
    \label{fig:09}
\end{figure}

\begin{figure}[H]
  \centering
  \includegraphics[width=0.9\textwidth]{10 - VoltageData with 2nd filter.png}
    \caption{\small نتائج تحليل \el{FFT} للجهد بعد اضافة المرشح الثاني - تظهر \el{THD} بنسبة $6.35\%$}
    \label{fig:10}
\end{figure}

\newpage

\begin{itemize}
  \item أظهرت النتائج في الصورة (الشكل \ref{fig:09}) المزيد من التحسن في نسبة تشوه التيار بعد اضافة المرشح الثاني، لتصل إلى \el{\boldmath $11.94\%$}.
  \item أظهرت النتائج في الصورة (الشكل \ref{fig:10}) تحسن آخر لا بأس به في نسبة تشوه الجهد واستقراره عند \el{\boldmath $6.35\%$}.
\end{itemize}

\vspace{20pt}

    \item \textbf{مرشح التوافقية الحادية عشرة ($n=11$):}\\[5pt] تم إدخال القيم: $R=0.115 \Omega$ و $L=0.0025 H$.
    
    \begin{figure}[H]
    \centering
    \includegraphics[width=0.8\linewidth]{12 - 3rd filter parameters.png} % Replace with actual file
    \caption{بارامترات مرشح التوافقية الحادية عشرة في ماتلاب}
    \label{fig:12}
\end{figure}

\newpage

    \begin{figure}[H]
      \vspace{-30pt}
        \centering
        \includegraphics[angle = 90, width=0.95\linewidth]{11 - Circuit with 3rd filter.png}
        \caption{شكل الدارة النهائي مع كافة المرشحات}
    \end{figure}
\end{enumerate}

\newpage
بإضافة المرشح الأخير المخصص للتوافقية الحادية عشرة ($550 Hz$)، نصل إلى الحالة الشبه مثالية المطلوبة في الوظيفة:
\begin{figure}[H]
  \centering
  \includegraphics[width=0.9\textwidth]{13 - CurrentData with 3rd filter.png}
    \caption{\small نتائج تحليل \el{FFT} للتيار بعد اضافة المرشح الثالث - تظهر \el{THD} بنسبة $9.93\%$}
    \label{fig:13}
\end{figure}

\begin{figure}[H]
  \centering
  \includegraphics[width=0.9\textwidth]{14 - VoltageData with 3rd filter.png}
    \caption{\small نتائج تحليل \el{FFT} للجهد بعد اضافة المرشح الثالث - تظهر \el{THD} بنسبة $4.76\%$}
    \label{fig:14}
\end{figure}

\newpage

\begin{itemize}
    \item \textbf{تحليل التيار:} تظهر الصورة (الشكل \ref{fig:13}) استمرار انخفاض نسبة التشوه الكلي للتيار واختفاء معظم المركبات التوافقية الضارة.
    \item \textbf{تحليل الجهد:} تظهر الصورة (الشكل \ref{fig:14}) انخفاض نسبة \el{$THD$} للجهد إلى أدنى مستوياتها، مما يجعل موجة الجهد جيبية ونقية بشكل كبير.
\end{itemize}
